\section{Two-box alternating tensor cores}
\label{sec:alternating-two-tail}

\paragraph{Status and scope.}
This complete new author derivation has passed whole-proof readings
by the root and independent /root/imprimitive reviewer, with separate
exact computational supplements. It concerns the actual tensor-core representations
with ordinary alternating labels $(n-2,2)$ and $(n-2,1,1)$, including
all compatible scalars and factor permutations. The inducing character
field may be larger than the original imaginary quadratic field.
Single factors retain the prior D-UP and D-EX exclusions; the new
source theorem concerns at least two factors.

For $n\ge15$, let $P_n$ be the permutation module on letters,
$U_n$ its deleted summand, and $P_{2,n}$ the permutation module
on unordered pairs. We use the following two rational orthogonal
modules, whose restrictions to $A_n$ are irreducible:
\[
\begin{array}{c|c|c|c}
W&\text{ordinary label}&a=\dim W&b\\\hline
W_{\mathrm{UP}}=P_{2,n}\ominus P_n&(n-2,2)&n(n-3)/2&1\\
W_{\mathrm{EX}}=\Lambda^2U_n&(n-2,1,1)&(n-1)(n-2)/2&\lfloor(n-1)/2\rfloor.
\end{array}                                                    \tag{ATT.1}
\]
The actual submodule and quotient realizations, including the
ordinary labels and irreducibility, are those in D-UP-01/03 and
D-EX-01/03. They are not virtual representations in what follows.

\begin{lemma}[D-ATT-01: exact minimum support and lower traces]
\label{D-ATT-01}
For every nonidentity $\pi\in S_n$,
\[
 \nu(W_{\mathrm{UP}}(\pi))\ge n-3,
 \qquad \nu(W_{\mathrm{EX}}(\pi))\ge n-2.                    \tag{ATT.2}
\]
Both bounds are attained by a transposition. Every permutation
has character at least $-b$ on the corresponding module in
\textup{(ATT.1)}.
\end{lemma}

\begin{proof}
Write $k=\operatorname{ind}(\pi)$ and
$D_\lambda=\operatorname{rank}(1-\lambda W(\pi))$.
The already proved all-order inequalities (UP.3)--(UP.4) give,
when $k\ge2$,
\[
 D_1\ge n-5+n/4\ge n-3,
 \qquad
 D_{\lambda\ne1}\ge n-5+(n-4)/2\ge n-3.                   \tag{ATT.3}
\]
Their hypotheses are precisely a nonidentity permutation in
the first inequality and a nonidentity target scalar in the
second; both hold here. If $k=1$, the permutation is a
transposition. Its negative eigenspace on $P_{2,n}$ has
dimension $n-2$, while that on $P_n$ has dimension one.
Thus the two eigenspace dimensions on $W_{\mathrm{UP}}$ are
$n-3$ and $a-(n-3)$, with the latter larger. This proves
the first exact minimum.

For the exterior square, (EX.4) gives
$D_1\ge n-4+n/3\ge n-2$ when $k\ge2$.
For every nonidentity target scalar, (EX.3) gives
\[
 D_{\lambda\ne1}\ge
 a-\lfloor(n-1)^2/4\rfloor\ge n-2.                       \tag{ATT.4}
\]
For the last inequality it suffices to omit the floor:
$a-(n-1)^2/4-(n-2)=(n^2-8n+11)/4>0$ for $n\ge15$.
A transposition on $U_n$ has a negative line and a positive
space of dimension $n-2$. Its exterior square has negative
space of dimension $n-2$ and a larger positive space, proving
the second exact minimum. Taking the minimum over all
$\lambda$ proves (ATT.2); targets which are not eigenvalues
have full moving rank and cause no exception.

Let $f_1$ and $f_2$ count the fixed letters and the two-cycles
of $\pi$. The actual permutation-module and exterior-square
character identities give
\[
 \chi_{\mathrm{UP}}(\pi)=\frac{f_1(f_1-3)}2+f_2,
 \qquad
 \chi_{\mathrm{EX}}(\pi)
     =\frac{(f_1-1)(f_1-2)}2-f_2.                          \tag{ATT.5}
\]
The first is at least $-1$: the only negative values of its
first summand occur at $f_1=1,2$, where it is $-1$.
For the second, if $f_1=0$ it is at least $1-\lfloor n/2\rfloor$;
if $f_1=1,2$ it is at least $-\lfloor(n-1)/2\rfloor$;
and if $f_1\ge3$, its first summand is nonnegative and
$f_2\le\lfloor(n-3)/2\rfloor$. This proves the stated floor
in every parity. The identities follow respectively by counting
fixed unordered pairs and from
$\chi_{\Lambda^2U}(g)=(\chi_U(g)^2-\chi_U(g^2))/2$,
where $\chi_U(g)=f_1-1$ and $\chi_U(g^2)=f_1+2f_2-1$.
\end{proof}

Put $V=W^{\otimes t}$, $e=a^t$, and let $\mathcal R$ be the
actual rational orthogonal tensor-wreath image of $S_n\wr S_t$.
Set
\[
 c=b/a,
 \qquad
 \mu=\begin{cases}2/n,&W=W_{\mathrm{UP}},\\
                  2/(n-1),&W=W_{\mathrm{EX}}.
       \end{cases}                                       \tag{ATT.6}
\]

\begin{lemma}[D-ATT-02: scalar extraction and marked inertias]
\label{D-ATT-02}
Let a finite linear group $J$ normalize the ordinary core $A_n^t$
on $V$. Then every $g\in J$ is uniquely $\lambda(g)T(g)$,
with $T(g)\in\mathcal R$ and a root of unity $\lambda(g)$.
The scalar coordinate is an actual character with image $\mu_M$,
and the actual character field $F$ of $J$ is contained in
$\mathbf Q(\zeta_M)$. Every genuine inertia has moving rank
at least $\mu e$.

For every $T\in\mathcal R$ and integer $r$,
\[
                      \operatorname{tr}(T^r)\ge-ce.        \tag{ATT.7}
\]
If $\lambda(g)\notin\mu_6$, its actual moving rank is at least
$(3-c)e/4$. In a generating product-one tuple, at least two
entries have this property whenever its original imaginary
quadratic field $K\subseteq F$ differs from $\mathbf Q(\sqrt{-3})$.
If $8\mid|\Delta_K|$, two scalar orders are divisible by eight
and those entries move at least $3e/4$. If an odd prime
$p\mid|\Delta_K|$, two scalar orders are divisible by $p$
and those entries move at least $(1-1/(p-1))e$.
\end{lemma}

\begin{proof}
The intrinsic simple factors and their matrix algebras are
permuted by every normalizer. Remove the corresponding tensor
permutation. Since $\operatorname{Aut}(A_n)=S_n$ for this range,
the remaining factor automorphisms are implemented by the
chosen actual $S_n$ operators on $W$. Removing them leaves a
scalar by Schur's lemma. This is the same full-automorphism
input and argument as D-AST-02, now on the two rational
irreducible modules in (ATT.1). Sign choices for extensions
to $S_n$ are absorbed in the scalar; no outer coset is omitted.

A scalar in $\mathcal R$ centralizes every represented simple
factor, so its factor permutation is trivial. Each of its
remaining $S_n$ permutations then centralizes $A_n$, since
the nontrivial $A_n$ representation is faithful. That
centralizer is trivial: a permutation commuting with every
three-cycle preserves each three-element support, and their
intersections fix every letter. Thus
$\mathcal R\cap\mathbf C^*I=\{I\}$. The decomposition is
unique; multiplication shows both coordinates are
homomorphisms, and finite order makes the scalar image
$\mu_M$. Rationality of $T$ gives
$\operatorname{tr}(g)\in\mathbf Q(\zeta_M)$ and hence the
asserted field containment. Equality $F=K$ is not used.

For a factor-preserving nonscalar tensor operator, the
largest eigenspace proportion is at most the largest such
proportion on any active factor, by convolution of its finite
spectral distributions. Lemma D-ATT-01 gives support at least
$\mu e$. For a nontrivial factor cycle of length $l\ge2$,
conjugate factorwise so that only the cycle product remains
as a twist, and diagonalize that finite-order product.
In the resulting tensor-word basis, each cyclic word orbit
contributes at most one vector to any chosen eigenspace.
There are $a$ constant words; all other word orbits have
length at least two. Thus the maximum is at most
$(a^l+a)/2\le3a^l/4$. The remaining tensor factors cannot
increase this proportion. Quarter support is stronger than
$\mu$ here. A nonidentity scalar has full moving rank,
which proves the assertion for every genuine inertia.

Trace contraction over the cycles of a tensor permutation
expresses $\operatorname{tr}(T)$ as a product of character
values of $W$ on their cycle products. Each value lies in
$[-b,a]$. If a product of $s\le t$ such values is negative,
one factor has absolute value at most $b$, and all others
at most $a$. Consequently it is at least
$-ba^{s-1}\ge-ba^{t-1}=-ce$. Every power belongs to the
same finite tensor-wreath group, proving (ATT.7).

For a root $\xi$ of order $d$, rationality of the actual
$T$ makes all primitive $d$th roots have the same multiplicity
$m_\xi$, so $\varphi(d)m_\xi\le e$. If $\xi\notin\mu_6$,
this gives $m_\xi\le e/4$ except possibly when $d=4$.
In that case the distinct $i,-i$ eigenspaces have equal
multiplicity $m_\xi$, and
\[
                 -ce\le\operatorname{tr}(T^2)
                                    \le e-4m_\xi.           \tag{ATT.8}
\]
Thus always $m_\xi\le(1+c)e/4$ for such a root. Apply this
to the actual $\lambda(g)^{-1}$ eigenspace of $T(g)$.
If the scalar order is divisible by eight or an odd prime
$p$, its totient is at least four or $p-1$, respectively.
These prove all marked-support bounds, including the case
that the requisite eigenvalue is absent.

For completeness, the arithmetic marks use the exact
argument of HD-AST-04, with all its representation premises
just verified here. If $K\ne\mathbf Q(\sqrt{-3})$, the
scalar image cannot be contained in $\mu_6$, since otherwise
$K\subseteq F\subseteq\mathbf Q(\zeta_6)$ would force equality.
The images of a generating product-one tuple in
$\mu_M/(\mu_M\cap\mu_6)$ cannot have exactly one nonidentity
entry. There are therefore at least two marked entries.

The standard quadratic conductor criterion gives
$|\Delta_K|\mid M$ from $K\subseteq\mathbf Q(\zeta_M)$.
We use Milne, Chapter V, Example3.2 p.~157, Remark3.8 p.~158,
and Examples3.10--3.11 p.~159, at exactly the primary-source
scope recorded for HD-AST-04.\cite{HD:MilneCFT2020}
The complete invoked passages were freshly read here.
For a prime $p\mid M$, the quotient
$\mu_M/\mu_{M/p}$ is nontrivial, and at least two tuple
entries have nonidentity images. If $p^v\Vert M$, each
such scalar has order divisible by $p^v$. This proves the
odd-prime and conductor-eight marks. If $M\equiv2\pmod4$,
the usual replacement of its cyclotomic conductor by $M/2$
does not affect divisibility by the original $M$.
\end{proof}

\begin{theorem}[D-ATT-03: exclusions for two-box tensor cores]
\label{D-ATT-03}
In the actual sixfold source setting, suppose the faithful
inducing image has ordinary core $A_n^t$ on either module
$W^{\otimes t}$ of (ATT.1), with $n\ge15$ and $t\ge2$.
Keep the original imaginary quadratic field only as
$K\subseteq F$. There is no period-image rank-nine source
in any of the following ranges:
\[
\begin{array}{c|c|c}
\text{field condition}&W_{\mathrm{UP}}&W_{\mathrm{EX}}\\\hline
K\ne\mathbf Q(\sqrt{-3})&15\le n\le39&15\le n\le39\\
8\mid|\Delta_K|&15\le n\le39&15\le n\le40\\
p\mid|\Delta_K|,\ p\ge7\text{ prime}
 &15\le n\le10(p-1)-1&15\le n\le10(p-1).
\end{array}                                                    \tag{ATT.9}
\]
In particular the remaining necessary lower bounds in the
eight non-Eisenstein class-number-one fields are
\[
\begin{array}{c|rrrrrrrr}
K&\mathbf Q(i)&\mathbf Q(\sqrt{-2})&\mathbf Q(\sqrt{-7})&
\mathbf Q(\sqrt{-11})&\mathbf Q(\sqrt{-19})&
\mathbf Q(\sqrt{-43})&\mathbf Q(\sqrt{-67})&\mathbf Q(\sqrt{-163})\\\hline
W_{\mathrm{UP}}:\ n\ge&40&40&60&100&180&420&660&1620\\
W_{\mathrm{EX}}:\ n\ge&40&41&61&101&181&421&661&1621.
\end{array}                                                    \tag{ATT.10}
\]
\end{theorem}

\begin{proof}
Use the actual Shapiro/Hurwitz period factorization established
in D-BI-01, not a subtraction of dimensions on the original
base. If $h$ is the intermediate genus and $D$ the sum of
actual moving ranks, multiplicity six gives
$6=2e(h-1)+D$. Here $e\ge a^2\ge8100$.
Genus at least two is impossible; genus one gives $D=6$,
at most six genuine branch marks and period-image rank at
most six. Rank nine therefore forces genus zero, at least
twelve genuine marks, and
\[
                              D=2e+6.                       \tag{ATT.11}
\]
The generating product-one inertia tuple is the actual one
of this intermediate faithful image.

For a non-Eisenstein original field, two marked entries from
D-ATT-02 and ten other genuine entries give
\[
  2e+6\ge\left(\frac{3-c}{2}+10\mu\right)e,
  \qquad \left(10\mu-\frac{1+c}{2}\right)e\le6.             \tag{ATT.12}
\]
For $W_{\mathrm{UP}}$, the coefficient times $a^2$ is
\[
 \frac{n(n-3)}8\big((40-n)(n-3)-2\big).                    \tag{ATT.13}
\]
On $15\le n\le39$, it is at least
$15\cdot12\cdot10/8>6$. Since the coefficient is positive
and $e\ge a^2$, this contradicts (ATT.12).
For $W_{\mathrm{EX}}$, use $c\le1/(n-2)$. The coefficient
times $a^2$ is at least
\[
 \frac{(n-1)(n-2)}8\big((40-n)(n-2)-1\big)
       \ge\frac{14\cdot13\cdot12}{8}>6                   \tag{ATT.14}
\]
on the same interval. This proves the first row.

For conductor divisible by eight, replace the two marked
supports by $3e/4$. The coefficient becomes $10\mu-1/2$.
The UP range already follows from the first row. In the EX
range $15\le n\le40$, its product with $a^2$ is
\[
              \frac{(n-1)(n-2)^2(41-n)}8>6,                \tag{ATT.15}
\]
so this also excludes the additional endpoint.

Finally put $h_0=p-1\ge6$. The two scalar orders divisible
by $p$ give the necessary inequality
$(10\mu-2/h_0)e\le6$. For UP its coefficient is positive
when $n<10h_0$, and its product with $a^2$ is
\[
             \frac{n(n-3)^2(10h_0-n)}{2h_0}.              \tag{ATT.16}
\]
On $15\le n\le10h_0-1$, concavity of $n(10h_0-n)$
puts its minimum at an endpoint, and that minimum is
$10h_0-1$: the other endpoint is $15(10h_0-15)$ and
their difference is $140h_0-224>0$.
Since $(n-3)^2\ge144$, expression (ATT.16) is at least
$144(10h_0-1)/(2h_0)>6$.
For EX put $x=n-1$, so $14\le x\le10h_0-1$.
The corresponding expression is
\[
             \frac{x(x-1)^2(10h_0-x)}{2h_0}.              \tag{ATT.17}
\]
Again the endpoint minimum of $x(10h_0-x)$ is $10h_0-1$,
as $14(10h_0-14)-(10h_0-1)=130h_0-195>0$.
Now $(x-1)^2\ge169$ proves the contradiction. This proves
all claimed ranges; (ATT.10) substitutes the exact quadratic
discriminants $4,8,7,11,19,43,67,163$.
\end{proof}

\paragraph{Strictness and remaining gate.}
For either label at $n=39$ and $t=2$, the previous ALT07
necessary support inequality is satisfied: its fraction is
$\beta(39,2)=36/703<1/10$. Thus the new exclusion over
$\mathbf Q(i)$ strictly narrows the previous main gate.
No claim is made that such a tuple or source existed before
the new obstruction. The conductor rows are compatible with
the separate broader long-label conductor refinement, but
use the sharper exact two-box support minima. All larger
unexcluded parameters and other long labels remain open in
this result; an all-degree tensor exclusion has not been
deduced from these inequalities.

The finite supplement checks exact cycle spectra, both
character formulas and support minima, tensor-cycle traces,
and the rational arithmetic. The all-degree proofs above,
complete compatible-normalizer extraction, original-field
containment and actual period-factorization interfaces are
not inferred from that finite calculation. The conditional
status of Markman, Severn and032 is unchanged, and no cycle
supplier is asserted by this obstruction.
